\documentclass[10pt,a4paper]{amsart}
\usepackage[margin=19mm]{geometry}
\usepackage{amsmath,amssymb,mathtools}
\usepackage[scr=rsfs,cal=euler]{mathalfa}
\usepackage{xfrac}
\usepackage[unicode,hidelinks,bookmarks=false]{hyperref}
\newcommand{\ie}{i.e.\ }
\newcommand{\cf}{cf.\ }
\newcommand{\resp}{respectively\ }
\newcommand{\Subsection}{\subsection}
\providecommand{\nobreakdash}{\nobreak\hbox{-}}
\setlength{\emergencystretch}{2em}
\multlinegap0pt
\usepackage{bm}
\usepackage{bbm}
\usepackage{dsfont}
\usepackage{xspace}
\usepackage{cancel}
\usepackage{mathtools}
\usepackage{amsthm}
\makeatletter
\@ifundefined{theorem}{\newtheorem{theorem}{Theorem}}{}
\@ifundefined{corollary}{\newtheorem{corollary}[theorem]{Corollary}}{}
\@ifundefined{lemma}{\newtheorem{lemma}[theorem]{Lemma}}{}
\@ifundefined{proposition}{\newtheorem{proposition}[theorem]{Proposition}}{}
\makeatother
\let\hat=\widehat
\newcommand {\Id}{\mathbb I}
\newcommand\C{{\mathbb C}}
\newcommand\R{{\mathbb R}}
\newcommand\Z{{\mathbb Z}}
\newcommand\cC{{\mathcal{C}}}
\newcommand\cL{{\mathcal{L}}}
\newcommand{\dist}{\mathrm{dist}}
\title[Stability Analysis of a Simple Discretization Method for a C]{Stability Analysis of a Simple Discretization Method for a Class of Strongly Singular Integral Equations}
\author{Martin Costabel, Monique Dauge, Khadijeh Nedaiasl}
\date{Source: arXiv:2302.13159v3 (CC BY 4.0)}
\begin{document}
\maketitle

\noindent\emph{Source and licence.} Stability Analysis of a Simple Discretization Method for a Class of Strongly Singular Integral Equations. Version arXiv:2302.13159v3 (CC BY 4.0), distributed under the Creative Commons Attribution 4.0 International licence, as recorded in the arXiv licence metadata for this exact version and shown on the abstract page https://arxiv.org/abs/2302.13159v3. Full rights notice: licenses/arxiv-2302.13159-CC-BY-4.0.txt.\par\smallskip

\noindent\textit{Editorial scope.} Counted unit: the Corollary whose source label is \texttt{C:stabEX2} in the article (source0.tex lines 1297--1303, source label \texttt{C:stabEX2}), delivered together with the article's own complete original proof of it (source0.tex lines 1305--1371, one \texttt{proof} environment of 67 lines). No mathematics is written, repaired or re-derived: statement and proof are the article's own text line for line. Required context retained uncounted: P:NumSym (lines 876--890); E:FFintsum (lines 965--976); E:FPintsum (lines 965--976); E:H0 (lines 986--995); E:H1 (lines 986--995); P:FInt (lines 1020--1036); L:heat (lines 1040--1052); E:DDA2dEx2 (lines 1243--1250); E:FEx2 (lines 1259--1266); L:im(Fex2) (lines 1279--1286); E:Fpos (lines 1282--1285). Every definition, display and label of those lines is retained unchanged. Omitted from this package and not redistributed: 1--875, 891--964, 977--985, 996--1019, 1037--1039, 1053--1242, 1251--1258, 1267--1278, 1286--1296, 1304--1304, 1372--1923. Nothing inside a retained excerpt is omitted or altered: every retained body is delivered exactly as the article's lines are written, so no omission receipt of any kind accompanies this package. Citations: the retained excerpts carry no citation command at all, so no bibliography entry of the article is transcribed and no reference is redistributed. Reference adaptation: none was needed and none was made; every reference inside the delivered unit resolves inside it, so no unresolved reference or citation is produced. Printed slips kept literal: no printed word, symbol, hypothesis or number of the counted statement or of its proof was altered. Layout and class: the article's own class is \texttt{amsart} with packages including mathalfa, amsmath, amsfonts, amsbsy, amssymb, ulem, times, mathrsfs, exscale, graphicx, hyperref, dsfont, color, booktabs; the wrapper is the lane's pinned \texttt{amsart} editorial wrapper at 10pt on A4, which restates the theorem-like environments the retained text needs and adds the article's own macro definitions that the retained text needs (\texttt{\textbackslash C}, \texttt{\textbackslash Id}, \texttt{\textbackslash R}, \texttt{\textbackslash Z}, \texttt{\textbackslash cC}, \texttt{\textbackslash cL}, \texttt{\textbackslash dist}, \texttt{\textbackslash hat}), with extra wrapper packages (bm, bbm, dsfont, xspace, cancel, mathtools). The delivered numbering is the unit's own. Assets: no retained span contains a figure environment, a graphics-inclusion command, a PDF-page inclusion, a table or a TikZ picture; no image file is copied into this package, the package declares no asset dependency and no image file is redistributed with it. Licence: version arXiv:2302.13159v3 of this article is distributed under the Creative Commons Attribution 4.0 International licence, as recorded in the arXiv licence metadata for this exact version and shown on the abstract page https://arxiv.org/abs/2302.13159v3. A full rights notice is delivered at licenses/arxiv-2302.13159-CC-BY-4.0.txt. Characters: every retained excerpt is pure ASCII, so the delivered engine, XeLaTeX with its default Latin Modern text font, reports no missing character.
\begin{proposition}
 \label{P:NumSym}
The symbol $F(\tau)$ of the infinite Toeplitz matrix $T$ is a bounded function given for any $\beta>0$ by the exponentially convergent sums
\begin{equation}
\label{E:Fewald}
 F(\tau) = 
 \sum_{m\in\Z^{d},m\ne0}K(m)\tfrac{\Gamma(\tfrac d2+1,\beta^{2}|m|^{2})}{\Gamma(\frac d2+1)}e^{im\cdot \tau} +
 \sum_{n\in\Z^{d}}\hat K(\tau+2\pi n)\,e^{-\frac{|\tau+2\pi n|^{2}}{4\beta^{2}}}\,.
\end{equation}
In the period cube $Q=[-\pi,\pi]^{d}$, it is $C^{\infty}$ outside of $\,0$, and it has the form
\begin{equation}
\label{E:F=F0+Khat} 
 F(\tau) = \hat K(\tau) + F_{0}(\tau) \quad\mbox{ where $F_{0}$ is analytic in $Q$ and $F_{0}(0)=0$}.  
\end{equation}
\end{proposition}

\begin{align}
\label{E:FFintsum}
 F^{F}(\xi) &= \int_{\beta^{2}}^{\infty}\!\!\! H^{F}(\xi,s)\,ds \quad\mbox{ with }\quad
 H^{F}(\xi,s) = \sum_{m\in\Z^{d}} 
    \frac{p(m)}{\Gamma(\frac d2+1)} s^{\frac d2}e^{-|m|^{2}s}
    \,e^{im\cdot \xi}\\
\label{E:FPintsum}
 F^{P}(\xi) &= \int_{0}^{\beta^{2}}\!\!\! H^{P}(\xi,s)\,ds \quad\mbox{ with }\quad
 H^{P}(\xi,s) =
 \sum_{n\in\Z^{d}}
    \frac{-\pi^{\frac d2}p(\xi+2\pi n)}{4\Gamma(\frac d2+1)} s^{-2}e^{-\frac{|\xi+2\pi n|^{2}}{4s}}\,.
\end{align}

\begin{align}
\label{E:H0} 
 H_{0}(\xi,s) &= -\tfrac{\pi^{\frac d2}}{4\Gamma(\frac d2+1)}
 \!\!\!
  \sum_{n\in\Z^{d},n\ne0} \!\!\!
    p(\xi+2\pi n) s^{-2}e^{-\frac{|\xi+2\pi n|^{2}}{4s}}\,,\\
\label{E:H1}
  H_{1}(\xi,s) &= -\tfrac{\pi^{\frac d2}}{4\Gamma(\frac d2+1)}
     p(\xi) s^{-2}e^{-\frac{|\xi|^{2}}{4s}}\,.
\end{align}

\begin{proposition}
 \label{P:FInt}
 The symbol $F(\xi)$ has the integral representation
\begin{equation}
\label{E:F=intH}
 F(\xi) = \int_{0}^{\infty} H(\xi,s)\,ds\,,
\end{equation}
where $H(\xi,s)$ is given either by the Fourier series $H^{F}$ in \eqref{E:FFintsum} or, equivalently, by the lattice sum $H^{P}$ in \eqref{E:FPintsum}.  
The decomposition $F=F_{0}+\hat K$ in Proposition~\ref{P:NumSym} corresponds to the decomposition $H=H_{0}+H_{1}$ with $H_{0}$ and $H_{1}$ defined in \eqref{E:H0} and \eqref{E:H1}, and there holds
\begin{equation}
\label{E:F01Int}
F_{0}(\xi) = \int_{0}^{\infty} H_{0}(\xi,s)\,ds \qquad
\mbox{ and }\quad
\hat K(\xi) = \int_{0}^{\infty} H_{1}(\xi,s)\,ds \,.
\end{equation}
In these integrals, the functions $s\mapsto H_{0}(\xi,s)$, $s\mapsto H_{1}(\xi,s)$, and $s\mapsto H(\xi,s)$ are integrable on $(0,\infty)$ for any $\xi\in Q$, for any $\xi\in\R^{d}\setminus\{0\}$, and for any $\xi\in Q\setminus\{0\}$, respectively.
\end{proposition}

\begin{lemma}
The functions
 \label{L:heat}
$$
 (\xi,s)\mapsto s^{-\frac d2}H_{0}(\xi,s),\quad
 (\xi,s)\mapsto s^{-\frac d2}H_{1}(\xi,s),\quad
 (\xi,s)\mapsto s^{-\frac d2}H(\xi,s),\quad
$$
are solutions of the heat equation 
$$
(\partial_{s}-\Delta_{\xi})u(\xi,s)=0\quad\mbox{ in }\; Q\times(0,\infty).
$$
\end{lemma}

\begin{equation}
\label{E:DDA2dEx2}
 \lambda u_{m} + \frac{1}{\pi} N^{-2}
 \sum_{n\in\omega^{N},m\ne n} 
 \frac{(x^{N}_{m,1}-x^{N}_{n,1})(x^{N}_{m,2}-x^{N}_{n,2})}{|x^{N}_{m}-x^{N}_{n}|^{4}}
  u_{n}
 = f(x^{N}_{m}) \,,\quad (m\in\omega^{N})\,.
\end{equation}

\begin{equation}
\label{E:FEx2}
 F(\tau) = -\sum_{m\in\Z^{2},m\ne0}\frac{m_{1}m_{2}}{\pi|m|^{4}}\,e^{im\cdot \tau}
  = \frac4\pi \sum_{m_{1},m_{2}=1}^{\infty}
    \frac{m_{1}m_{2}}{(m_{1}^{2}+m_{2}^{2})^{2}}
    \sin(m_{1}\tau_{1}) \sin(m_{2}\tau_{2})
        \,.
\end{equation}

\begin{lemma}
 \label{L:im(Fex2)}
 Let $F(\tau)$ be as defined in \eqref{E:FEx2}. Then for any $\tau\in Q$
\begin{equation}
\label{E:Fpos}
  F(\tau) \in \cC=[-\frac12,\frac12].
\end{equation}
\end{lemma}

\begin{corollary}
 \label{C:stabEX2}
 Let $\cC=[-\frac12,\frac12]$ and $\lambda\in\C\setminus\cC$. Then for any $N$ the linear system~\eqref{E:DDA2dEx2} has a unique solution, and there is a uniform resolvent estimate \begin{equation}
\label{E:ResTNex2}
 \|(\lambda\Id-T^{N})^{-1}\|_{\cL(\ell^{2}(\Omega^{N}))} \le \dist(\lambda,\cC)^{-1}\,.
\end{equation}
\end{corollary}

\begin{proof}
For symmetry reasons, it is sufficient to prove \eqref{E:Fpos} for $\tau\in Q_{++}=(0,\pi)^{2}$.
For the proof of Lemma~\ref{L:im(Fex2)}, we will show the following:
\begin{equation}
\label{E:FposF0neg}
\mbox{ For any }\xi\in Q_{++}\,,\quad
F(\xi)\ge0 \;\mbox{ and } \; F_{0}(\xi)\le0\,.
\end{equation}
This implies $0\le F(\xi)\le \hat K(\xi)\le\frac12$, hence \eqref{E:Fpos}\,.

We use the integral representations from Proposition~\ref{P:FInt}
\begin{equation}
\label{E:F0Int}
 F_{0}(\xi) = \int_{0}^{\infty} H_{0}(\xi,s)ds\,,\quad
 \hat K(\xi) = \int_{0}^{\infty} H_{1}(\xi,s)\,ds\,,\quad
 F(\xi) = \int_{0}^{\infty} H(\xi,s)\,ds\,.
\end{equation}
According to \eqref{E:H0}, $H_{0}(\xi,s)=H(\xi,s)-H_{1}(\xi,s)$ with
\begin{equation}
\label{E:H0H1Ex2}
 H_{1}(\xi,s)= \frac{\xi_{1}\xi_{2}}{4s^{2}} e^{-\frac{|\xi|^{2}}{4s}}\,,\quad
 H(\xi,s)=  \sum_{n\in\Z^{2}} H_{1}(\xi+2\pi n,s)\,.
\end{equation}
Let $0\le\epsilon<T$ and $\Sigma_{\epsilon}^{T}=Q_{++}\times (\epsilon,T)$.
In $\Sigma_{\epsilon}^{T}$, we want to use the maximum principle for the heat equation (see Lemma~\ref{L:heat}) for the functions 
$\tilde H_{0}(\xi,s)=s^{-1}H_{0}(\xi,s)$ and $\tilde H(\xi,s)=s^{-1}H(\xi,s)$. 

Since $\tilde H_{1}(\xi,s)=s^{-1}H_{1}(\xi,s)$ is continuous for $(\xi,s)\in\R^{2}\times[0,\infty)\setminus\{0,0\}$ and the Poisson series 
$$
 \tilde H_{0}(\xi,s) =  \sum_{n\in\Z^{2},n\ne0} \tilde H_{1}(\xi+2\pi n,s)
$$
converges uniformly for $(\xi,s)\in\overline{\Sigma_{0}^{T}}$ for all $T>0$, we see that 
$\tilde H_{0}$ is continuous in $\overline{\Sigma_{0}^{T}}$ with initial value 
$\tilde H_{0}(\xi,0)=0$. 
On the lateral boundary we use the Fourier representation (see \eqref{E:FFintsum})
$$
 \tilde H(\xi,s) = \frac4\pi \sum_{m_{1},m_{2}=1}^{\infty}
    m_{1}m_{2} e^{-|m|^{2}s}
    \sin (m_{1}\xi_{1}) \sin(m_{2}\xi_{2})\,.
$$
If $\xi_{1}$ or $\xi_{2}$ is in $\{0,\pi\}$, this implies that $\tilde H=0$ and therefore
$$
 \tilde H_{0}(\xi,s) = -\tilde H_{1}(\xi,s)\le0 
 \quad\mbox{ for }(\xi,s)\in\partial Q_{++}\times(0,T]\,.
$$
According to Lemma \ref{L:heat}, $\tilde H_{0}$ satisfies the heat equation
$(\partial_{s}-\Delta_{\xi})\tilde H_{0}=0$ in $\Sigma_{0}^{T}$.
Thus we can apply the maximum principle to $\tilde H_{0}$ and obtain 
$\tilde H_{0}(\xi,s)\le0$ in $\Sigma_{0}^{T}$, hence also $H_{0}(\xi,s)\le0$.
Integrating over $s\in(0,\infty)$ yields
$$
 F_{0}(\xi) \le 0 \quad\mbox{ for }\xi\in Q_{++}\,.
$$

For $\tilde H$, we cannot apply the maximum principle directly in $\Sigma_{0}^{T}$, because $\tilde H$ is not continuous at $(0,0)\in\overline{\Sigma_{0}^{T}}$, but we can apply it in
$\Sigma_{\epsilon}^{T}$ for any $0<\epsilon<T$. On the lateral boundary, $\tilde H$ vanishes as seen above, and for the initial value at $s=\epsilon$ we have
$$
 \tilde H(\xi,\epsilon) = \tilde H_{0}(\xi,\epsilon) + \tilde H_{1}(\xi,\epsilon)
 \ge \tilde H_{0}(\xi,\epsilon) \ge \delta(\epsilon)
$$
with $\delta(\epsilon)=\inf_{\xi\in Q_{++}}\tilde H_{0}(\xi,\epsilon)$.
Hence by the maximum principle, in $\overline{\Sigma_{\epsilon}^{T}}$ we have 
$$
 \tilde H(\xi,s) \ge \min\{0,\delta(\epsilon)\}\,.
$$
Now, as we have seen above, $\tilde H_{0}(\cdot,s)$ tends to $0$ uniformly as $s\to0$, hence $\delta(\epsilon)\to0$ as $\epsilon\to0$, which implies $\tilde H(\xi,s)\ge0$ for any $s>0$ and $\xi\in Q_{++}$. After integrating over $s$, we finally get $F(\xi)\ge0$ for $\xi\in Q_{++}$, and the proof of the Lemma is complete.
\end{proof}

\end{document}
